%% kapa: exam=1 round=27 year=2016 period=1 form=A subject=economics subject_ko=경제학원론 seq=32 number=72
%% source: Q-Net 감정평가사 기출문제 (cst00309) artlSeq=5169402 file=extracted/1008/5169402/제27회 감정평가사 1차 문제/제27회 감정평가사 1교시문제(A형).pdf
%% conversion: kordoc 4.24.51 -> tools/gen1.py (auto); review=proofread-sonnet; flags=none
\begin{question}{72}
어떤 경제의 총수요곡선과 총공급곡선이 각각 $P = -Y^D + 2$, $P = P^e + (Y^S - 1)$이다. $P^e$가 1.5일 때, 다음 설명 중 옳은 것을 모두 고른 것은? (단, $P$는 물가수준, $Y^D$는 총수요, $Y^S$는 총공급, $P^e$는 기대물가수준이다.)
\begin{framedtext}
ㄱ. 이 경제의 균형은 $P = 1.25$, $Y = 0.75$이다. \par
ㄴ. 이 경제는 장기 균형 상태이다. \par
ㄷ. 합리적 기대 가설 하에서는 기대물가수준 $P^e$는 1.25이다. \par
\end{framedtext}
\begin{choices}
  \item ㄱ
  \item ㄴ
  \item ㄱ, ㄷ
  \item ㄴ, ㄷ
  \item ㄱ, ㄴ, ㄷ
\end{choices}
\end{question}
\answer{1}
